% Pipeline Audebert et al. (2020), digambar ulang dari Gambar 1 dan Bagian 3-4 paper.
% Dipakai bersama oleh laporan dan presentasi; butuh TikZ (positioning, arrows.meta).
\resizebox{\linewidth}{!}{%
\begin{tikzpicture}[
  node distance=4mm and 7mm,
  kotak/.style={draw, rounded corners=2pt, align=center, font=\small,
                minimum height=10mm, inner sep=3pt},
  citra/.style={kotak, fill=blue!10},
  teks/.style={kotak, fill=orange!15},
  fusi/.style={kotak, fill=green!15},
  panah/.style={-{Stealth[length=2mm]}, thick}]
  \node[kotak] (in) {Citra dokumen\\(grayscale)};
  % alur citra
  \node[citra, above right=4mm and 9mm of in] (rs) {Resize 384$\times$384\\gray$\to$3 kanal};
  \node[citra, right=of rs] (mb) {MobileNetV2\\(ImageNet, fine-tune)};
  \node[citra, right=of mb] (gap) {GAP 1280-d\\$\to$ FC 128-d};
  % alur teks
  \node[teks, below right=4mm and 9mm of in] (ocr) {Tesseract 4.0\\(LSTM, EN)};
  \node[teks, right=of ocr] (ft) {SpaCy token\\+ FastText (Common Crawl)};
  \node[teks, right=of ft] (cnn) {CNN 1D 4 lapis\\$\to$ 128-d};
  % fusi
  \node[fusi, right=12mm of $(gap)!0.5!(cnn)$] (cat) {Konkatenasi\\256-d};
  \node[fusi, right=of cat] (out) {MLP +\\softmax};
  \draw[panah] (in) |- (rs);  \draw[panah] (rs) -- (mb);  \draw[panah] (mb) -- (gap);
  \draw[panah] (in) |- (ocr); \draw[panah] (ocr) -- (ft); \draw[panah] (ft) -- (cnn);
  \draw[panah] (gap) -| (cat); \draw[panah] (cnn) -| (cat); \draw[panah] (cat) -- (out);
  \node[font=\footnotesize\itshape, below=1mm of out, align=center] {dilatih\\end-to-end};
\end{tikzpicture}}
